\section*{Chapitre 18 - Aires}
\addcontentsline{toc}{section}{Chapitre 18 - Aires}


\subsection*{I. Aire}
\addcontentsline{toc}{subsection}{I. Aire}
\begin{Définition}[c]
L'\textbf{aire} d'une figure est la mesure de sa surface.
\end{Définition}

\begin{Exemple}
\begin{minipage}{0.45\textwidth}
Donner l'aire des deux figures ci-contre. 

\ligne

\ligne

\end{minipage} \hfill
\begin{minipage}{0.55\textwidth}
\begin{center}
\begin{tikzpicture}[baseline,scale = 0.5]

    \tikzset{
      point/.style={
        thick,
        draw,
        cross out,
        inner sep=0pt,
        minimum width=5pt,
        minimum height=5pt,
      },
    }

    	\draw[color ={gray},opacity = 0.4] (0,0)--(0,4);
	\draw[color ={gray},opacity = 0.4] (1,0)--(1,4);
	\draw[color ={gray},opacity = 0.4] (2,0)--(2,4);
	\draw[color ={gray},opacity = 0.4] (3,0)--(3,4);
	\draw[color ={gray},opacity = 0.4] (4,0)--(4,4);
	\draw[color ={gray},opacity = 0.4] (5,0)--(5,4);
	\draw[color ={gray},opacity = 0.4] (6,0)--(6,4);
	\draw[color ={gray},opacity = 0.4] (7,0)--(7,4);
	\draw[color ={gray},opacity = 0.4] (8,0)--(8,4);
	\draw[color ={gray},opacity = 0.4] (9,0)--(9,4);
	\draw[color ={gray},opacity = 0.4] (10,0)--(10,4);
	\draw[color ={gray},opacity = 0.4] (11,0)--(11,4);
	\draw[color ={gray},opacity = 0.4] (12,0)--(12,4);
	\draw[color ={gray},opacity = 0.4] (13,0)--(13,4);
	\draw[color ={gray},opacity = 0.4] (14,0)--(14,4);
	\draw[color ={gray},opacity = 0.4] (15,0)--(15,4);
	\draw[color ={gray},opacity = 0.4] (0,0)--(15,0);
	\draw[color ={gray},opacity = 0.4] (0,1)--(15,1);
	\draw[color ={gray},opacity = 0.4] (0,2)--(15,2);
	\draw[color ={gray},opacity = 0.4] (0,3)--(15,3);
	\draw[color ={gray},opacity = 0.4] (0,4)--(15,4);
	\draw[color={Couleur8},line width = 1.5,preaction={fill,color = {Couleur8!25}, opacity = 0.5}] (1,0)--(1,1)--(0,1)--(0,2)--(2,2)--(2,3)--(3,3)--(3,2)--(4,2)--(4,1)--(5,1)--(5,0)--(2,0)--cycle;
	\draw[color={Couleur8},line width = 1.5,preaction={fill,color = {Couleur8!25}, opacity = 0.5}] (7,1)--(7,3)--(6,3)--(6,4)--(8,4)--(8,2)--(10,2)--(10,0)--(9,0)--(9,1)--(8,1)--cycle;
	\draw [color={black}] (3,0.5) node[anchor = center,scale=1, rotate = 0] {{\small Figure 1}};
	\draw [color={black}] (6.9,1.5) node[anchor = west,scale=1, rotate = 0] {{\small Figure 2}};

	\draw[color={Couleur8},preaction={fill,color = {Couleur8!25}, opacity = 0.5}] (13,2)--(14,2)--(14,3)--(13,3)--cycle;
	\draw (13.5,1.25) node[anchor = center, rotate=0] {\normalsize \color{black}{$1$ unité }};
		\draw (13.5,0.5) node[anchor = center, rotate=0] {\normalsize \color{black}{d'aire}};

\end{tikzpicture}
\end{center}
\end{minipage}

\end{Exemple}


\vspace{-0.2cm}
\subsection*{II. Unités d'aire}
\addcontentsline{toc}{subsection}{II. Unités d'aire}

\begin{Règle}
L'unité d'aire usuelle est le mètre carré (noté m²) qui représente l'aire d'un carré d'un mètre de côté. On utilise aussi : ses sous-multiples (dm², cm², mm²).
\end{Règle}

\begin{center}
\begin{tikzpicture}

 % Troisième ligne - noms des positions
\node[draw, Couleur1,fill=Couleur1!20, minimum width=2.3cm, minimum height=0.9cm, align=center] at (-2,0.85) {km²};
\node[draw, Couleur1, fill=Couleur1!20,minimum width=2.3cm, minimum height=0.9cm, align=center] at (0.3,0.85) {hm²};
\node[draw, Couleur1, fill=Couleur1!20,minimum width=2.3cm, minimum height=0.9cm, align=center] at (2.6,0.85) {dam²};
\node[draw, Couleur4, fill=Couleur4!20,minimum width=2.3cm, minimum height=0.9cm, align=center] at (4.9,0.85) {m²};
\node[draw,Couleur6, fill=Couleur6!20,minimum width=2.3cm, minimum height=0.9cm, align=center] at (7.2,0.85) {dm²};
\node[draw, Couleur6, fill=Couleur6!20,minimum width=2.3cm, minimum height=0.9cm, align=center] at (9.5,0.85) {cm²};
\node[draw, Couleur6, fill=Couleur6!20,minimum width=2.3cm, minimum height=0.9cm, align=center] at (11.8,0.85) {mm²};

% Quatrième ligne - les chiffres du nombre (2 colonnes par unité)
% km² - 2 colonnes
\node[draw, Couleur1, minimum width=1.15cm, minimum height=0.7cm] at (-2.575,0.05) {};
\node[draw, Couleur1, minimum width=1.15cm, minimum height=0.7cm] at (-1.425,0.05) {};

% hm² - 2 colonnes  
\node[draw, Couleur1, minimum width=1.15cm, minimum height=0.7cm] at (-0.275,0.05) {};
\node[draw, Couleur1, minimum width=1.15cm, minimum height=0.7cm] at (0.875,0.05) {};

% dam² - 2 colonnes
\node[draw, Couleur1, minimum width=1.15cm, minimum height=0.7cm] at (2.025,0.05) {};
\node[draw, Couleur1, minimum width=1.15cm, minimum height=0.7cm] at (3.175,0.05) {};

% m² - 2 colonnes
\node[draw, Couleur4, minimum width=1.15cm, minimum height=0.7cm] at (4.325,0.05) {};
\node[draw, Couleur4, minimum width=1.15cm, minimum height=0.7cm] at (5.475,0.05) {};

% dm² - 2 colonnes
\node[draw, Couleur6, minimum width=1.15cm, minimum height=0.7cm] at (6.625,0.05) {};
\node[draw, Couleur6, minimum width=1.15cm, minimum height=0.7cm] at (7.775,0.05) {};

% cm² - 2 colonnes
\node[draw, Couleur6, minimum width=1.15cm, minimum height=0.7cm] at (8.925,0.05) {};
\node[draw, Couleur6, minimum width=1.15cm, minimum height=0.7cm] at (10.075,0.05) {};

% mm² - 2 colonnes
\node[draw, Couleur6, minimum width=1.15cm, minimum height=0.7cm] at (11.225,0.05) {};
\node[draw, Couleur6, minimum width=1.15cm, minimum height=0.7cm] at (12.375,0.05) {};

% Cinquième ligne - les chiffres du nombre (2 colonnes par unité)
% km² - 2 colonnes
\node[draw, Couleur1, minimum width=1.15cm, minimum height=0.7cm] at (-2.575,-0.65) {};
\node[draw, Couleur1, minimum width=1.15cm, minimum height=0.7cm] at (-1.425,-0.65) {};

% hm² - 2 colonnes  
\node[draw, Couleur1, minimum width=1.15cm, minimum height=0.7cm] at (-0.275,-0.65) {};
\node[draw, Couleur1, minimum width=1.15cm, minimum height=0.7cm] at (0.875,-0.65) {};

% dam² - 2 colonnes
\node[draw, Couleur1, minimum width=1.15cm, minimum height=0.7cm] at (2.025,-0.65) {};
\node[draw, Couleur1, minimum width=1.15cm, minimum height=0.7cm] at (3.175,-0.65) {};

% m² - 2 colonnes
\node[draw, Couleur4, minimum width=1.15cm, minimum height=0.7cm] at (4.325,-0.65) {};
\node[draw, Couleur4, minimum width=1.15cm, minimum height=0.7cm] at (5.475,-0.65) {};

% dm² - 2 colonnes
\node[draw, Couleur6, minimum width=1.15cm, minimum height=0.7cm] at (6.625,-0.65) {};
\node[draw, Couleur6, minimum width=1.15cm, minimum height=0.7cm] at (7.775,-0.65) {};

% cm² - 2 colonnes
\node[draw, Couleur6, minimum width=1.15cm, minimum height=0.7cm] at (8.925,-0.65) {};
\node[draw, Couleur6, minimum width=1.15cm, minimum height=0.7cm] at (10.075,-0.65) {};

% mm² - 2 colonnes
\node[draw, Couleur6, minimum width=1.15cm, minimum height=0.7cm] at (11.225,-0.65) {};
\node[draw, Couleur6, minimum width=1.15cm, minimum height=0.7cm] at (12.375,-0.65) {};
  
\end{tikzpicture}
\end{center}

\begin{Exemples}
\vspace{-0.6cm}
\begin{multicols}{2}
\begin{itemize}[label=\textcolor{Couleur8}{$\bullet$}]
\item $54,2\text{ m}^2 = $ \ligne

\item $7,5\text{ cm}^2 = $ \ligne

\end{itemize}
\end{multicols}
\vspace{-0.6cm}
\end{Exemples}




\begin{Propriété}[c] \textbf{Formulaire.}
\begin{itemize}[label=\textcolor{Couleur1}{$\bullet$}]
\item L'aire $A$ d'un rectangle de longueur $L$ et de largeur $l$ est donné par la formule :\\
 $A_{\text{Rectangle}} = L \times l$.

\item L'aire $A$ d'un carré de côté $c$ est donné par la formule :  $A_{\text{Carré}}=c \times c$

\end{itemize}
\end{Propriété}




\begin{Exemples}
\vspace{-0.8cm}
\begin{multicols}{2}
\begin{minipage}{0.15\textwidth}
\vspace{-2cm}
\begin{align*}
\textcolor{Couleur8}{\bullet} A_{\text{LKJI}}&=\ldots\ldots\\
&= \ldots\ldots\\
&= \ldots\ldots
\end{align*}
\end{minipage} \hfill
\begin{minipage}{0.35\textwidth}
\begin{center}
\begin{tikzpicture}[baseline,scale = 0.3]

    \tikzset{
      point/.style={
        thick,
        draw,
        cross out,
        inner sep=0pt,
        minimum width=5pt,
        minimum height=5pt,
      },
    }
    \clip (-2,-2.14) rectangle (10.14,10);
    	\draw [color={black}] (-0.5,-0.5) node[anchor = center,scale=1, rotate = 0] {I};
	\draw [color={black}] (8.5,-0.64) node[anchor = center,scale=1, rotate = 0] {J};
	\draw [color={black}] (8.64,8.36) node[anchor = center,scale=1, rotate = 0] {K};
	\draw [color={black}] (-0.36,8.5) node[anchor = center,scale=1, rotate = 0] {L};
	\draw[color={Couleur6!75}] (8,-0.14)--(7.5,-0.13)--(7.51,0.37)--(8.01,0.36)--cycle;
	\draw[color={Couleur6!75}] (0.14,8)--(0.13,7.5)--(0.63,7.49)--(0.64,7.99)--cycle;
	\draw[color={Couleur6!75}] (8.14,7.86)--(7.64,7.87)--(7.63,7.37)--(8.13,7.36)--cycle;
	\draw[color={Couleur6!75}] (0,0)--(0.5,-0.01)--(0.51,0.49)--(0.01,0.5)--cycle;
	\draw [color={Couleur6!75}] (4,-0.07) node[anchor = center,scale=0.5, rotate = -1] {//};
\draw [color={Couleur6!75}] (8.07,3.86) node[anchor = center,scale=0.5, rotate = 89] {//};
\draw [color={Couleur6!75}] (4.14,7.93) node[anchor = center,scale=0.5, rotate = -1] {//};
\draw [color={Couleur6!75}] (0.07,4) node[anchor = center,scale=0.5, rotate = -271] {//};

	\draw [color={black}] (3.99,-0.77) node[anchor = center,scale=1, rotate = -1.0025699999999915] {8 cm};
	\draw[color={Couleur6},line width = 1.5] (0,0)--(8,-0.14)--(8.14,7.86)--(0.14,8)--cycle;

\end{tikzpicture}
\end{center}
\end{minipage} 


\vspace{-2cm}
\begin{minipage}{0.15\textwidth}
\begin{align*}
\textcolor{Couleur8}{\bullet}A_{\text{VUTS}}&= \ldots\ldots \\
&= \ldots\ldots\\
&= \ldots\ldots
\end{align*}
\end{minipage} \hfill
\begin{minipage}{0.35\textwidth}
\begin{center}
\begin{tikzpicture}[baseline,scale = 0.3]

    \tikzset{
      point/.style={
        thick,
        draw,
        cross out,
        inner sep=0pt,
        minimum width=5pt,
        minimum height=5pt,
      },
    }

\clip (-2.18,-2) rectangle (11.99,7.35);
    	\draw [color={black}] (-0.5,-0.5) node[anchor = center,scale=1, rotate = 0] {S};
	\draw [color={black}] (10.49,-0.15) node[anchor = center,scale=1, rotate = 0] {T};
	\draw [color={black}] (10.31,5.85) node[anchor = center,scale=1, rotate = 0] {U};
	\draw [color={black}] (-0.68,5.5) node[anchor = center,scale=1, rotate = 0] {V};
	\draw[color={Couleur6!50}] (9.99,0.35)--(9.49,0.33)--(9.47,0.83)--(9.97,0.85)--cycle;
	\draw[color={Couleur6!50}] (-0.18,5)--(-0.16,4.5)--(0.34,4.52)--(0.32,5.02)--cycle;
	\draw[color={Couleur6!50}] (9.81,5.35)--(9.31,5.33)--(9.33,4.83)--(9.83,4.85)--cycle;
	\draw[color={Couleur6!50}] (0,0)--(0.5,0.02)--(0.48,0.52)--(-0.02,0.5)--cycle;
	%\draw [color={gray}] (5,0.18) node[anchor = center,scale=1, rotate = 2] {//};
\draw [color={Couleur6!50}] (5,0.18) node[anchor = center,scale=0.5, rotate = -358] {//};

%	\draw [color={gray}] (4.82,5.18) node[anchor = center,scale=1, rotate = -358] {//};
\draw [color={Couleur6!50}] (4.82,5.18) node[anchor = center,scale=0.5, rotate = 2] {//};

%	\draw [color={gray}] (9.9,2.85) node[anchor = center,scale=1, rotate = -88] {/};
\draw [color={Couleur6!50}] (9.9,2.85) node[anchor = center,scale=0.5, rotate = -88] {/};

%	\draw [color={gray}] (-0.09,2.5) node[anchor = center,scale=1, rotate = -88] {/};
\draw [color={Couleur6!50}] (-0.09,2.5) node[anchor = center,scale=0.5, rotate = -88] {/};

	\draw [color={black}] (5.02,-0.52) node[anchor = center,scale=1, rotate = -357.99346] {10 cm};
	\draw [color={black}] (10.6,2.88) node[anchor = center,scale=1, rotate = -87.93824] {5 cm};
	\draw[color={Couleur6!75},line width = 1.5] (0,0)--(9.99,0.35)--(9.81,5.35)--(-0.18,5)--cycle;

\end{tikzpicture}
\end{center}
\end{minipage} 

\end{multicols}


\begin{minipage}{0.6\textwidth}
\textcolor{Couleur8}{$\bullet$} Un triangle rectangle est la moitié d'un rectangle donc : 
\begin{align*}
A_{\text{PNO}}&=\ldots\ldots\ldots\ldots \\
&= \ldots\ldots\ldots\ldots\\
&= \ldots\ldots\ldots\ldots \\
&= \ldots\ldots\ldots\ldots
\end{align*}
\end{minipage} \hfill
\begin{minipage}{0.4\textwidth}
\begin{center}
\begin{tikzpicture}[baseline,scale = 0.4]

    \tikzset{
      point/.style={
        thick,
        draw,
        cross out,
        inner sep=0pt,
        minimum width=5pt,
        minimum height=5pt,
      },
    }
    \clip (-2,-2) rectangle (9.2,5);
    	\draw [color={black}] (-0.5,-0.5) node[anchor = center,scale=1, rotate = 0] {N};
	\draw [color={black}] (7.7,-0.25) node[anchor = center,scale=1, rotate = 0] {O};
	\draw [color={black}] (-0.1,3.5) node[anchor = center,scale=1, rotate = 0] {P};
	\draw[color={Couleur6!25}] (0,0)--(0.8,0.03)--(0.77,0.83)--(-0.03,0.8)--cycle;
	\draw[color={Couleur6!50},line width = 1.5] (0,0)--(7.2,0.25)--(-0.1,3)--cycle;
	\draw [color={black}] (3.9,2.57) node[anchor = center,scale=1, rotate = -20.64201] {{\small 26 cm}};
	\draw [color={black}] (3.63,-0.87) node[anchor = center,scale=1, rotate = -358.01135999999997] {{\small 24 cm}};
	\draw [color={black}] (-1.05,1.47) node[anchor = center,scale=1, rotate = -88.09085] {{\small 10 cm}};

\end{tikzpicture}
\end{center}
\end{minipage} 

\end{Exemples}

